% sec:gr-efe — The Einstein Field Equations
\section{The Einstein Field Equations}
\label{sec:gr-efe}

The geometric machinery of \cref{ch:differential-geometry} culminated in a
remarkable accident: the Einstein tensor $\einstein$ is divergence-free purely
as a consequence of the Riemannian geometry, with no physics assumed
(\cref{eq:contracted-bianchi}).  Meanwhile, any sensible theory of matter
demands that the stress-energy tensor $\stressenergy$ --- which encodes the
density, pressure, and momentum flux of whatever fills spacetime --- satisfy
the conservation law $\nabla^\mu \stressenergy = 0$.  Two symmetric,
divergence-free $(0,2)$ tensors built from entirely different starting
points: one from the geometry, the other from the matter.  Einstein's stroke
was to set them proportional to each other.
\physnote{The requirement $\nabla^\mu \stressenergy = 0$ generalises the
flat-spacetime conservation laws $\pd{\mu} T^{\mu\nu} = 0$ to curved
spacetime.  On a curved manifold, covariant divergence replaces ordinary
divergence.}

\paragraph{The Newtonian correspondence.}

Any candidate field equation for gravity must pass a basic sanity check: in
the limit of weak fields and slow-moving matter, it must reduce to the
Poisson equation that encodes Newtonian gravity.  This requirement pins down
the proportionality constant.  Newton's gravitational potential $\Phi$
satisfies
%
\begin{equation}
  \laplacian \Phi = 4\pi\rho \,,
  \label{eq:poisson}
\end{equation}
%
where $\rho$ is the mass density and we work in geometrised units
($G = c = 1$).  The left-hand side involves second spatial derivatives of the
potential; in general relativity, second derivatives of the metric play the
role of $\laplacian \Phi$, and the Ricci tensor is the natural candidate
since it is built from second derivatives of $\metric$
(\cref{eq:riemann-components}).  On the right-hand side, the mass density
$\rho$ is just one component of the stress-energy tensor --- specifically
$T_{00}$ in the rest frame of the matter.  The simplest generally covariant
equation with the correct Newtonian limit is
%
\begin{equation}
  \einstein = 8\pi\,\stressenergy \,,
  \label{eq:efe}
\end{equation}
%
where the factor of $8\pi$ is fixed by requiring that \cref{eq:efe}
reproduce \cref{eq:poisson} in the weak-field, low-velocity limit.
\Cref{sec:gr-weak-field} carries out this matching explicitly.
\histnote{Einstein arrived at \cref{eq:efe} in November 1915, after
several years of false starts.  The trace-reversed form
(\cref{eq:efe-trace-reversed} below) appeared first; the compact
$\einstein = 8\pi\stressenergy$ followed once the Einstein tensor and the
contracted Bianchi identity were recognised as a package.}

Expanding the Einstein tensor using its definition
(\cref{eq:einstein-tensor}), the field equations read
%
\begin{equation}
  \ricci - \tfrac{1}{2}\,\metric\,\ricciscalar = 8\pi\,\stressenergy \,.
  \label{eq:efe-expanded}
\end{equation}
%
This is a system of ten equations (one for each independent component of the
symmetric tensors $\ricci$ and $\stressenergy$) for the ten independent
components of the metric $\metric$.  The count matches, but --- as we discuss
below --- four of the ten equations are constraints rather than evolution
equations, a subtlety tied to the coordinate freedom of general relativity.

\paragraph{Trace-reversed form.}

The trace of \cref{eq:efe-expanded} isolates the Ricci scalar in terms of
the matter content alone, which lets us rewrite the field equations in a more
convenient form.  Contracting with $\inversemetric$ and using
$\inversemetric\,\metric = \delta^\mu{}_\mu = 4$ (in four spacetime
dimensions) gives
%
\begin{equation}
  \ricciscalar - 2\ricciscalar = 8\pi\,T \,,
  \label{eq:efe-trace}
\end{equation}
%
where $T = \inversemetric\,\stressenergy$ is the trace of the stress-energy
tensor and the factor of $2$ arises as $\tfrac{1}{2} \times 4$ from the
second term.  Solving for the Ricci scalar, $\ricciscalar = -8\pi\,T$, and
substituting back into \cref{eq:efe-expanded} yields the trace-reversed form
%
\begin{equation}
  \ricci = 8\pi\!\left(\stressenergy
    - \tfrac{1}{2}\,\metric\,T\right) .
  \label{eq:efe-trace-reversed}
\end{equation}
%
The trace-reversed form is often more convenient for deriving solutions,
because the left-hand side involves only the Ricci tensor rather than the
Ricci tensor minus a trace term.  In vacuum, where $\stressenergy = 0$ and
hence $T = 0$, the trace equation gives $\ricciscalar = 0$, and
\cref{eq:efe-expanded} reduces to
%
\begin{equation}
  \ricci = 0 \,.
  \label{eq:vacuum-efe}
\end{equation}
%
Every exact spacetime in the codebase --- Schwarzschild, Kerr, and
Morris--Thorne --- is either a vacuum solution or satisfies a particularly
simple stress-energy tensor.  The vacuum equations $\ricci = 0$ are the
starting point for the Schwarzschild derivation in \cref{ch:schwarzschild}.
\xrefnote{The stress-energy tensors for perfect fluids and electromagnetic
fields are defined in \cref{sec:gr-stress-energy}.}

\paragraph{Structure as partial differential equations.}

The field equations look algebraic as written, but the Ricci tensor hides
second derivatives of the metric.  Since $\ricci$ is a contraction of the
Riemann tensor (\cref{eq:ricci-def}), and the Riemann tensor involves
derivatives of the Christoffel symbols (\cref{eq:riemann-components}), which
are themselves first derivatives of the metric
(\cref{eq:christoffel-formula}), the Einstein equations are a system of ten
coupled, second-order, nonlinear partial differential equations for the ten
components of $\metric$.  The nonlinearity enters through the
$\Gamma\Gamma$ terms in the Riemann tensor: schematically,
$R \sim \partial^2 g + (\partial g)^2$.  This is in sharp contrast to
electromagnetism, where Maxwell's equations are linear in the field.
The nonlinearity has a direct physical consequence: gravitational fields
gravitate.  The energy stored in the gravitational field itself contributes
to the curvature, so superposition fails and exact solutions are rare.
\physnote{The self-interaction encoded in the $(\partial g)^2$ terms is what
makes numerical solutions essential: even two-body problems in general
relativity have no closed-form solution.}

In a given coordinate system, the ten metric components
$g_{\mu\nu}(x^\alpha)$ are the unknowns, and the ten equations
(\cref{eq:efe-expanded}) relate them to whatever matter is present.  For
generic matter distributions, no closed-form solution exists, and the
equations must be solved numerically --- the task of the numerical-relativity
chapters in Part~IV.

\paragraph{Gauge freedom and the constraint count.}

General relativity is a diffeomorphism-invariant theory: the physics is
unchanged by smooth reparametrisations of the coordinates.  A coordinate
transformation $x^\mu \to x'^\mu(x)$ introduces four arbitrary functions
(one per spacetime dimension), so metrics that differ by a coordinate
transformation describe the same physical spacetime.  This gauge freedom
means that four of the ten metric components are pure gauge --- they can be
set to any convenient value by a coordinate choice, without affecting the
physics.

The contracted Bianchi identity is what enforces this counting at the level
of the equations themselves.  The four identities
$\nabla^\mu \einstein = 0$ (\cref{eq:contracted-bianchi}) are satisfied by
the left-hand side of \cref{eq:efe} regardless of the metric, so only six
of the ten field equations are truly independent.  In the initial-value
formulation, four of these take the form of constraints that the initial data
must satisfy, while the remaining six govern time evolution --- a structure
we revisit in detail in the ADM decomposition of
\cref{ch:adm-decomposition}.  The six independent equations match the six
physical degrees of freedom of the gravitational field (ten metric components
minus four gauge choices).
\warnnote{The gauge freedom of general relativity is more subtle than
electromagnetic gauge freedom.  In electromagnetism, gauge transformations
act on the potential but leave the field strength unchanged.  In general
relativity, the ``field'' {\itshape is} the metric, and gauge
transformations change its components directly.  Only geometric invariants
--- scalar curvature, geodesic lengths, etc.\ --- are physically
meaningful.}

\paragraph{Conservation from geometry.}

The contracted Bianchi identity has a second, equally important
consequence.  Since the left-hand side of \cref{eq:efe} is identically
divergence-free, the right-hand side must be as well wherever the field
equations hold:
%
\begin{equation}
  \nabla^\mu \stressenergy = 0 \,.
  \label{eq:stress-energy-conservation}
\end{equation}
%
Energy-momentum conservation is not an additional postulate appended to the
field equations --- it is a mathematical consequence of them.  Any
stress-energy tensor that satisfies \cref{eq:efe} automatically satisfies
\cref{eq:stress-energy-conservation}.  This is why the Einstein tensor,
rather than the Ricci tensor alone, appears on the left-hand side: the
combination $\ricci - \tfrac{1}{2}\metric\ricciscalar$ is the unique
divergence-free symmetric $(0,2)$ tensor built from the metric and its first
two derivatives that is linear in the second derivatives (up to an additive
$\Lambda\,\metric$ term, discussed below).
\histnote{Hilbert derived the field equations from a variational principle
(the Einstein--Hilbert action
$S = \int \ricciscalar\,\detmetric\,\vol{4}$) almost simultaneously with
Einstein in November 1915.  The variational approach makes the Bianchi
identity and conservation law manifest as consequences of diffeomorphism
invariance, via Noether's second theorem.}

\paragraph{The cosmological constant.}

The Einstein tensor is not the only divergence-free symmetric tensor that
can appear on the left-hand side.  The metric itself is covariantly constant
($\cd{\alpha}\,\metric = 0$ by metric compatibility), so
$\Lambda\,\metric$ is divergence-free for any constant $\Lambda$.  The most
general field equation consistent with the requirements above is therefore
%
\begin{equation}
  \einstein + \Lambda\,\metric = 8\pi\,\stressenergy \,,
  \label{eq:efe-lambda}
\end{equation}
%
where $\Lambda$ is the cosmological constant.  Cosmological observations
constrain $\Lambda \approx 1.1 \times 10^{-52}\;\text{m}^{-2}$ (in SI
units), a scale so much larger than any black hole that $\Lambda$ has no
observable effect on the spacetimes we simulate.  We set $\Lambda = 0$ for
the remainder of this book.
\physnote{For a solar-mass black hole, $M_\odot \approx 1.5 \times
10^3\;\text{m}$ in geometrised units, so $\Lambda M^2 \sim 10^{-46}$ --- negligible.  Even for a supermassive hole with $M \sim 10^9 M_\odot$, the
ratio is of order $10^{-28}$.  The cosmological constant becomes relevant
only at cosmological distances (hundreds of megaparsecs), where it drives the
accelerating expansion of the universe.}

\paragraph{The vacuum equations and the road ahead.}

The spacetimes of primary interest in this book --- isolated black holes and
binary mergers --- are vacuum or nearly vacuum solutions.  Outside a black
hole's event horizon, $\stressenergy = 0$, and the field equations reduce to
$\ricci = 0$.  These ten equations (of which six are independent) are the
starting point for the exact solutions of
\cref{ch:schwarzschild,ch:kerr} and the numerical evolutions of Part~IV.
The next section introduces the stress-energy tensor for the matter sources
that do appear --- perfect fluids in accretion discs and electromagnetic
fields in jets --- before \cref{sec:gr-weak-field} verifies that the full
equations reduce to Newtonian gravity where they should.
